\chapter*{模块四：统计与概率}
\addcontentsline{toc}{chapter}{模块四：统计与概率}


\begin{quote}
数学不仅帮助我们精确计算，还帮助我们在\textbf{不确定}的世界中做出合理判断。统计与概率正是处理不确定性的数学工具。
\end{quote}


\section*{主线脉络}


统计与概率围绕一个核心问题展开：\textbf{如何用数据说话？}

在低年级，我们从身边最简单的活动开始——数一数、分一分、画一画，学习把杂乱的信息变成清晰的图表。随着年级升高，我们不再满足于"看图说话"，而是追问更深层的问题：这组数据的"中心"在哪里？数据的波动有多大？如果只调查一部分人，能推断全体的情况吗？一件事发生的可能性有多大？

这条主线可以概括为：

\[
\text{收集数据} \longrightarrow \text{整理描述} \longrightarrow \text{分析推断} \longrightarrow \text{把握不确定性}
\]


\section*{各章简介}


\begin{center}
\begin{tabular}{llll}
\toprule
章节 & 标题 & 适用年级 & 核心内容 \\
\midrule
§4.1 & 数据的收集与整理 & 1-6 年级 & 分类计数、统计表、条形图、折线图、扇形图、数据收集方法 \\
§4.2 & 平均数、中位数与众数 & 5-7 年级 & 用一个数来代表一组数据的"中心"位置 \\
§4.3 & 数据的波动程度 & 8 年级 & 极差、方差——衡量数据的离散程度 \\
§4.4 & 频率分布直方图与频率分布表 & 7-8 年级 & 把数据分组，用频率刻画数据的分布 \\
§4.5 & 概率初步 & 7-9 年级 & 随机事件、等可能事件的概率、列表法与树形图 \\
§4.6 & 总体与样本 & 8-9 年级 & 抽样调查、用样本估计总体 \\
\bottomrule
\end{tabular}
\end{center}


\section*{学习建议}


统计与概率和其他模块的关系非常紧密：计算平均数离不开四则运算（§1.2），方差的计算用到平方和分数（§1.5、§1.6），概率常以百分数的形式出现（§1.7）。学习本模块时，请确保对应的"数与运算"基础已经扎实。

与纯粹的计算不同，统计问题往往\textbf{没有唯一的标准答案}——同一组数据，用平均数还是中位数来描述更合适？取决于具体的情境。在学习过程中，请养成"先思考再计算"的习惯：\textbf{为什么选这个统计量？它告诉了我们什么？}
